\pnum{5.29}{Three particles distributed among three orbitals}
There are three distinguishable particles and
three available orthonormal one-body states.
Each labeled particle independently has
three choices, so
$\boxed{\mathcal N_D=3^3=27}$.
For three identical bosons, an arrangement
is specified by nonnegative occupancies
$n_a+n_b+n_c=3$. Stars-and-bars gives
\begin{equation}
\boxed{\mathcal N_B=
\binom{3+3-1}{3}=\binom53=10}.
\end{equation}
For three identical fermions with no additional
spin degree of freedom, the Pauli principle
allows at most one particle in each orbital,
so all three must be singly occupied:
\begin{equation}
\boxed{\mathcal N_F=\binom33=1}.
\end{equation}
The last state is the unique antisymmetric
Slater determinant up to phase.

\pnum{5.30}{Fermi energy in a two-dimensional square well}
For a large two-dimensional well, count
states in momentum space in the
thermodynamic limit. With area $A$,
each $k$-state occupies $(2\pi)^2/A$
and each accepts two spin orientations.
A filled disk of radius $k_F$ has
\begin{equation}
N_e=2\frac{A}{(2\pi)^2}\pi k_F^2
=\frac{Ak_F^2}{2\pi}.
\end{equation}
Let $\sigma=N_e/A$ be the areal density.
Then
$\boxed{k_F=\sqrt{2\pi\sigma}}$,
and
\begin{equation}
\boxed{E_F=\frac{\sh^2k_F^2}{2m}
=\frac{\pi\sh^2\sigma}{m}}.
\end{equation}
Boundary effects vanish relative to
the area in the macroscopic limit.

\pnum{5.31}{Cohesive-energy estimate including spin in three dimensions}
In the textbook's idealized box analogy,
take $N$ isolated monovalent atoms, each
hosting an electron in a cubic box of
side $a$. The individual ground-state
energy is
\begin{equation}
E_{\rm atom}=
\frac{\pi^2\sh^2}{2ma^2}(1^2+1^2+1^2)
=\frac{3\pi^2\sh^2}{2ma^2}.
\end{equation}
The separated total is $NE_{\rm atom}$.
After combining into a cubic metal
of volume $V=Na^3$, the same $N$
electrons occupy free-electron states
with twofold spin degeneracy.
Their number density is
$n=N/V=1/a^3$, giving
\begin{equation}
E_F=\frac{\sh^2}{2ma^2}(3\pi^2)^{2/3},
\quad
E_{\rm metal}=\frac35 NE_F.
\end{equation}
Define the cohesive-energy gain
as the separated energy minus the
metal energy. Per atom,
\begin{equation}
\boxed{\frac{\Delta E_{\rm coh}}N=
\frac{\sh^2}{2ma^2}\left[
3\pi^2-\frac35(3\pi^2)^{2/3}\right]}.
\end{equation}
It is positive in this model, indicating
an energetic advantage to delocalizing
electrons across a much larger volume.
The comparison intentionally ignores
the detailed ionic binding potentials
and is only an illustrative toy model.

\pnum{5.32}{Fermi gas with unequal spin populations}
Let $N_\pm$ be the spin-labeled counts
in volume $V$ and $n_\pm=N_\pm/V$.
Each spin species has degeneracy one,
so its Fermi wavevector is
$k_{F,\pm}=(6\pi^2n_\pm)^{1/3}$.
Integrating the kinetic energy for
each species gives
\begin{equation}
\boxed{E_{\rm tot}=
\frac{3\sh^2}{10m}(6\pi^2)^{2/3}
V(n_+^{5/3}+n_-^{5/3})}.
\end{equation}
For equal populations $n_\pm=\rho/2$,
the result reduces to
$E_{\rm tot}/V=
(3/10)(\sh^2/m)(3\pi^2)^{2/3}
\rho^{5/3}$, the unpolarized formula.
With magnetization
$M=-\mu_B(n_+-n_-)$,
define $\eta=M/(\rho\mu_B)$, so
$n_\pm=\rho(1\mp\eta)/2$.
Substitute:
\begin{equation}
\frac{E_{\rm tot}}V=
\frac{3}{10}\frac{\sh^2}{m}
(3\pi^2)^{2/3}\rho^{5/3}
\frac{(1+\eta)^{5/3}+
(1-\eta)^{5/3}}2.
\end{equation}
The binomial expansion yields
$\frac12[(1+\eta)^{5/3}+
(1-\eta)^{5/3}]=
1+\frac59\eta^2+O(\eta^4)$.
Thus
\begin{equation}
\boxed{\frac{E_{\rm tot}}V=
\frac{\sh^2(3\pi^2\rho)^{5/3}}
{10\pi^2m}
\left[1+\frac59
\left(\frac{M}{\rho\mu_B}\right)^2
+O(M^4)\right]}.
\end{equation}
The positive quadratic term makes
$M=0$ the kinetic-energy minimum.

\pnum{5.33}{Pauli paramagnetism and aluminum susceptibility}
For a weak uniform field along $z$,
the Zeeman energy density is $-MB$.
From Problem 5.32, the kinetic
energy density for small $M$ is
\begin{equation}
\frac{E_{\rm kin}}V
=\frac35\rho E_F+
\frac{E_F}{3\rho\mu_B^2}M^2+O(M^4).
\end{equation}
Minimize
$\mathcal E(M)=E_{\rm kin}/V-MB$:
\begin{equation}
\frac{\partial\mathcal E}{\partial M}
=\frac{2E_F}{3\rho\mu_B^2}M-B=0,
\qquad
\boxed{M=\frac{3\rho\mu_B^2}{2E_F}B}.
\end{equation}
With susceptibility defined by
$\chi=\mu_0\dd M/\dd B$,
\begin{equation}
\boxed{\chi_P=
\frac{3\mu_0\rho\mu_B^2}{2E_F}},
\quad
E_F=\frac{\sh^2}{2m_e}(3\pi^2\rho)^{2/3}.
\end{equation}
For aluminum $\rho=18.1\times10^{22}
\,{\rm cm^{-3}}=1.81\times10^{29}\,{\rm m^{-3}}$,
substitution gives
$\boxed{\chi_P\simeq1.2\times10^{-5}}$.
The textbook's comparison value is
$\chi_{\rm exp}=22\times10^{-6}$.
This is spin-only Pauli susceptibility;
orbital diamagnetism and interactions
are intentionally omitted.

\pnum{5.34}{Stoner criterion in a pair-counting interaction model}
The model assigns extra energy $U$ to
each opposite-spin pair:
\begin{equation}
\Delta E=UN_+N_-.
\end{equation}
Since $N_++N_-=\rho V$ and
$N_+-N_-=-MV/\mu_B$,
\begin{equation}
\boxed{\Delta E=
\frac{UV^2}{4}
\left(\rho^2-\frac{M^2}{\mu_B^2}\right)}.
\end{equation}
From Problem 5.32 the kinetic
contribution quadratic in $M$ is
$V E_F M^2/(3\rho\mu_B^2)$.
Thus the total quadratic coefficient is
\begin{equation}
\frac{V}{\mu_B^2}
\left(\frac{E_F}{3\rho}
-\frac{UV}{4}\right).
\end{equation}
The unpolarized state ceases to be a local
minimum when this becomes negative, giving
\begin{equation}
\boxed{U>U_c=\frac{4E_F}{3\rho V}
=\frac{4E_F}{3(N_++N_-)}}.
\end{equation}
\checkit{$U$ is the energy per
opposite-spin \emph{pair}. For an
extensive thermodynamic interaction,
one normally scales $U\propto1/V$.
This simple instability condition
is the stated textbook model,
not a complete theory of ferromagnetism.}
